\begin{trapbox}
	\begin{description}[style=nextline, leftmargin=2.8em, labelsep=0.5em, itemsep=0.35em]
		\item[\textbf{\textcolor{CTrap}{易错点一（忽略绝对值）}}] 直接写出 $d=\dfrac{\vec{AP}\cdot\mathbf{n}}{|\mathbf{n}|}$，没有加绝对值符号。
		\item[\textbf{\textcolor{CTrap}{正确理解（距离非负）：}}] 距离是非负数，必须取绝对值：$d=\dfrac{|\vec{AP}\cdot\mathbf{n}|}{|\mathbf{n}|}$。当点积为负数时，取绝对值保证正数。
		\item[\textbf{\textcolor{CTrap}{错因分析（方向误解）：}}] 学生误以为 $\vec{AP}$ 与 $\mathbf{n}$ 一定同向，忽略了 $P$ 在平面两侧的可能性。
	\end{description}

	\medskip
	\begin{description}[style=nextline, leftmargin=2.8em, labelsep=0.5em, itemsep=0.35em]
		\item[\textbf{\textcolor{CTrap}{易错点二（法向量零向量）}}] 使用零向量作为法向量代入公式。
		\item[\textbf{\textcolor{CTrap}{正确理解（必须非零）：}}] 法向量必须是非零向量，否则分母为零且法向量无意义。在应用中应先确认法向量不是零向量。
		\item[\textbf{\textcolor{CTrap}{错因分析（计算失误或忽略前提）：}}] 求法向量时可能得到零向量（如三点共线），此时平面未确定，不能使用公式。
	\end{description}

	\medskip
	\begin{description}[style=nextline, leftmargin=2.8em, labelsep=0.5em, itemsep=0.35em]
		\item[\textbf{\textcolor{CTrap}{易错点三（误用向量投影公式）}}] 认为 $d=|\vec{AP}|\cos\langle\vec{AP},\mathbf{n}\rangle$，而忽略了将 $\mathbf{n}$ 单位化。
		\item[\textbf{\textcolor{CTrap}{正确理解（公式正确形式）：}}] 正确的投影长度应为 $|\vec{AP}|\dfrac{|\vec{AP}\cdot\mathbf{n}|}{|\vec{AP}||\mathbf{n}|}= \dfrac{|\vec{AP}\cdot\mathbf{n}|}{|\mathbf{n}|}$。必须保留分母 $|\mathbf{n}|$。
		\item[\textbf{\textcolor{CTrap}{错因分析（混淆投影概念）：}}] 误将普通投影公式（带单位方向向量）与点积直接对应，忘记法向量不是单位向量。
	\end{description}
\end{trapbox}
